\chapter{图示与公式速查}\label{app:figures}
\index{附录}

\section{TikZ 标注图}

下面是一张带坐标和标注的虚构图，用来说明“图源单独文件、正文引用、图注和坐标标签”如何配合。

\begin{figure}[htbp]
  \centering
  \begin{tikzpicture}
  \begin{axis}[
    width=0.78\linewidth,
    height=6cm,
    xlabel={时间 / \si{\second}},
    ylabel={温度 / \si{\celsius}},
    xmin=0, xmax=6,
    ymin=18, ymax=26,
    grid=major,
    grid style={draw=gray!20},
    legend style={draw=none, font=\small},
  ]
    \addplot[bookteal, thick, mark=*] coordinates {
      (0,20.1) (1,20.8) (2,21.9) (3,23.2) (4,24.1) (5,24.6)
    };
    \addlegendentry{示例测量}
    \addplot[gray, dashed, thick, domain=0:5.5] {0.85*x + 19.7};
    \addlegendentry{参考线}
    \node[font=\small, align=left, anchor=west] at (axis cs:3.4,19.1) {示例标注：\\趋势仅供排版};
  \end{axis}
\end{tikzpicture}

  \caption{示例测量点与参考线}\label{fig:annotated}
\end{figure}

\section{常用命令}

\begin{itemize}[leftmargin=1.6em]
  \item 单位：\code{\textbackslash qty\{12.5\}\{\textbackslash milli\textbackslash metre\}}
  \item 交叉引用：\code{\textbackslash cref\{fig:process\}}
  \item 索引：\code{\textbackslash index\{术语\}}
  \item 文献：\code{\textbackslash cite\{example2026\}}
\end{itemize}

\section{图示检查}

\begin{enumerate}[leftmargin=1.8em]
  \item 箭头方向和连接关系是否与文字一致。
  \item 标签在打印尺寸下是否可读。
  \item 图注是否说明“这张图在讲什么”。
  \item 正文是否至少引用图一次，例如\cref{fig:annotated}。
\end{enumerate}
